\section{Definability of compact polyhedra}

\cite{Blackburn_Rijke_Venema_2001} in chapter 3 gives us a definability result for classes of finite transitive frames

\begin{theorem}

A class of finite transitive frame is modally definable if and only if it's closed under finite disjoint union,
bounded morphic images and generated subframes.
\end{theorem}
\begin{definition}[Modally definable class of simplicial complexes]\hfill 

    Suppose we have a class of simplicial complexes, We say that it is modally definable if and only if the corresponding class of posets is.
\end{definition}

We want to translate these class conditions in terms of simplicial complexes and then realize them in terms of polyhedra.

The picture we want to have in mind when building classes is the following: 

\[\begin{tikzcd}[ampersand replacement=\&]
	\& C \\
	{[C]} \& \Scplx \& {|C|} \\
	\Pos \&\& \Pol
	\arrow[from=1-2, to=2-1]
	\arrow[hook, from=1-2, to=2-2]
	\arrow[from=1-2, to=2-3]
	\arrow[hook, from=2-1, to=3-1]
	\arrow["{[-]}", from=2-2, to=3-1]
	\arrow["{|-|}", from=2-2, to=3-3]
	\arrow[hook, from=2-3, to=3-3]
	\arrow["\Nn", shift left=3, from=3-1, to=2-2]
\end{tikzcd}\]

Keep in mind that all kripke frames that arise in this context are posets, in particular they are transitive.

\subsection{Disjoint union}

Let us start with disjoint unions: these are the easiest construction since they are the same in all three settings

\begin{proposition}
    For a pair of simplicial complexes $\Sigma, \Sigma'$, 
    $[\Sigma\sqcup\Sigma']=[\Sigma]\sqcup[\Sigma']$
\end{proposition}
\begin{proof}
    Let $\sigma$ be a simplex of $[\Sigma\sqcup\Sigma']$, then -as the union is disjoint- it must be a simplex of either $[\Sigma]$ or $[\Sigma']$, and vice versa.
\end{proof}

Similarly
\begin{proposition}
    For a pair of simplicial complexes $\Sigma,\Sigma'$, $|\Sigma\sqcup \Sigma'|=|\Sigma|\sqcup|\Sigma'|$
\end{proposition}
\begin{proof}
    There are direct ways to show it but we appeal to abstract nonsense and say that the geometric realization is the coend 
    \[\int^{n\in \N}\Delta^n\cdot \Sigma_n\] 
    thus the constructions commutes with colimits.

    This is not as clear, but theoretically better than a direct proof in $\R^n$ 
    since it eventually will allow us to work with arbitrarily-sized polyhedra too.
\end{proof}

This means that
\begin{lemma}
    A class of simplicial complex $C$ is closed under disjoint union if and only if its corresponding class of posets $[C]$ is closed under disjoint union
    if and only if its corresponding class of polyhedra $|C|$ is closed under disjoint union.
\end{lemma}

\subsection{Generated subframes}

Now suppose we have a simplicial complex $\Sigma$: What corresponds to generated subframes of $[\Sigma]$?

Recall that
\begin{definition}[Generated subframe]

    A \textbf{Generated subframe} of a kripke frame $(W,R)$, is a subset $W'\subseteq W$ closed upwards under $R$.
\end{definition}

Recall that the posets we are considering are the 'coface' posets, i.e. the posets whose 
elements are the faces of $\Sigma$ and whose relation is the opposite of the inclusion.

This means that a generated subframe is $\uparrow-$closed under $\supseteq$, thus $\downarrow-$closed under $\subseteq$.

In other words 
\begin{lemma}\hfill

    Let $C$ be a class of simplicial complexes, then $[C]$ is closed under generated subframes if and only if $C$ is closed under
    subcomplexes. 
\end{lemma}
\begin{proof}\hfill 
    Let $\Sigma$ be a simplicial complex and $\Sigma'$ a subcomplex. 
    As $\Sigma'$ is a complex as well, then it's downward-closed.
    
    Conversely, if $S'\leq[\Sigma]$ is a $\downarrow$-closed subposet under $\subseteq$, 
    then by definition it's closed under taking subsets, thus it's a simplicial complex.
\end{proof}

\begin{lemma}
    Let $C$ be a class of simplicial complexes, then $|C|$ is closed under subpolyhedra if and only if 
    $C$ is closed under subcomplexes and subdivisions.
\end{lemma}
\begin{proof}
    If $\Sigma'\leq \Sigma$ is a subcomplex then $|\Sigma'|\subseteq |\Sigma|$ is a closed subpolyhedron.

    Conversely, if $S\subseteq |\Sigma|$ is a closed subpolyhedron, then there exists a triangulation $\Sigma'$ of $|\Sigma|$
    triangulating $S$, i.e. with a subcomplex $\Sigma''\subseteq \Sigma'$ such that $|\Sigma''|=S$ (the latter is a known fact and a 
    proof can be found in \cite{DBLP} as well as \cite{PLTopology} or other combinatorial topology textbooks.)
\end{proof}

\subsection{bounded morphic images}

Lastly we need to interpret bounded morphism in our triad. Recall 
\begin{definition}\hfill 

    Let $(W,R)$ and $(W',R')$ be kripke frames and $f:W\to W'$ a function.

    $f$ is a \textbf{bounded morphism} if and only if 
    \begin{itemize}
        \item $wRv$ implies $f(w)R'f(v)$
        \item $f(w)R' v'$ implies that there exists $v$ such that $wRv$ and $f(v)=v'$.
    \end{itemize}
\end{definition}

Interpreting this as the coface relation this means 
\begin{itemize}
    \item $\sigma\subseteq \tau$ implies $f(\sigma)\subseteq f(\tau)$
    \item $\sigma'\subseteq f(\tau)$ implies that there exists $\sigma\subseteq \tau$ such that $f(\sigma)=\sigma'$
\end{itemize}

Let's define a class of simplicial maps:
\begin{definition}[Bisimplicial map]\hfill 
    
    Let $(V,\Sigma),(V',\Sigma')$ be simplicial complexes. $f:V\to V'$ is called a \textbf{bisimplicial map} if 
    \[\forall s\in \Pp(V), f(s)\in \Sigma' \iff s\in \Sigma\]
\end{definition}

\begin{lemma}\hfill 
    
    A simplicial map $f:(V,\Sigma)\to(V',\Sigma')$ induces a bounded morphism $[\Sigma]\to[\Sigma']$ if and only if it's bisimplicial.
\end{lemma}
\begin{proof}
    Suppose $f$ is a bisimplicial map, then suppose 
    $\sigma\subseteq \tau$, then $f(\sigma)$ and $f(\tau)$ are simplices and $f(\sigma)\subseteq f(\tau)$ as it's an image map.

    Now suppose $\sigma'=\subseteq f(\tau)$. Suppose $\tau=v_1,...,v_n$, thus $f(\tau)=f(v_1),...,f(v_n)$. Up to reordering, $\sigma'=f(v_1),...,f(v_k)$.
    As $f$ is bisimplicial, $f(v_1),...,f(v_k)=f(v_1,...,v_k)$ is a simplex, thus $\sigma:=v_1,...,v_k$ is a simplex and - by structure- a subsimplex of $v_1,...,v_n$. 
    
    Now conversely, suppose that $f$ is a bounded morphism $[\Sigma]\to[\Sigma']$, the (co)induced $f:\Sigma\to\Sigma'$ is a bisimplicial map.

    First we need to check that image of a vertex is a vertex. This is because vertices are simplices, and minimal ones at that:
    suppose we have a vertex $v$ such that $f(v)=\sigma'$. Then if $\sigma'$ is not a vertex, there exists $w'\subsetneq \sigma'=f(v)$.

    This means that there must exist $w\subseteq v$ such that $f(w)=w'$, but such $w$ does not exist as $v$ is minimal. $\to\gets$.

    Image of a simplex is a simplex as the map's domain and codomain are simplices. Suppose that  
    $f(v_1,...,v_n)$ is a simplex of $\Sigma'$ but $v_1,...,v_n$ is not a simplex of $\Sigma$. 
    
    applying $f$ to $v_1,...,v_n$ we obtain that $f(v_1,...,v_n)=f(v_1,...,v_n)$ (thus in particular $\subseteq$).
    As $f$ is a bounded morphism defined on simplices and $f(v_1,...,v_n)$ is an element of the image,
    there must exist a simplex $v_1,...,v_k\subseteq v_1,...,v_n$ such that $f(v_1,...,v_k)=f(v_1,...,v_n)$. 
    Then, since they must have the same cardinality, it must be the case that
    $v_1,...,v_k=v_1,...,v_n$, thus $v_1,...,v_n$ is a simplex of $\Sigma$.
\end{proof}

Now to interpret bisimplicial maps as maps of polyhedra: 

Suppose we have the map $|f|:|\Sigma|\to|\Sigma'|$ that is the realization of a map $f:\Sigma\to\Sigma'$. 

If $f$ is bisimplicial, then this results in the conditions
\begin{itemize}
    \item if $\Hull(v_1,...,v_k)\subseteq |\Sigma|$, then $f(\Hull(v_1,...,v_k))=\Hull(f(v_1),...,f(v_k))\subseteq |\Sigma'|$, and 
    \item if $\Hull(f(v_1),...,f(v_k))\subseteq |\Sigma'|$, then $\Hull(v_1,...,v_k)\subseteq |\Sigma|$
\end{itemize} 

This is equivalent to it being a retraction onto the image:
\begin{theorem}\hfill

    Suppose $f$ if a map $\Sigma\to\Sigma'$, then it's a bisimplicial map if and only if $f:\Sigma\to\im(f)$ is a retraction.
\end{theorem}
\begin{proof}\hfill 

    Suppose that $f$ is bisimplicial.

    First, note that $f$ is surjective onto the vertices of $\im(f)$. 
    Moreover if $\sigma'=v_1',...,v_k'$ is a simplex in $\im(f)$, then there exist $v_1,...,v_k$ such that $f(v_i)=v_i'$ and $v_1,...,v_k$ is a simplex in $\Sigma$.

    This means that $f:\Sigma\to\im f$ is epimorphic at the level of simplices as well. 
    Now let $g:\im(f)\to \Sigma$ be defined as $g(v_i')=v_i$ for some $v_i$ such that $f(v_i)=v_i'$

    This is a well-defined map, sending vertices to vertices, and $f(g(v_1',...,v_n'))=f(v_1,...,v_n)=v1',...,v_n'$ thus $fg=1_{\im(f)}$

    Conversely, suppose that $f$ is a section, i.e. there exists a $g:\im(f)\to \Sigma$ such that $fg=1_{\im(f)}$
    Then suppose $f(v_1),...,f(v_n)=v_1',...,v_n'$ is a simplex. then $g(v_1'),...,g(v_n')=v_1,...,v_n$ is a simplex since $g$ is simplicial.
\end{proof}

\begin{corollary}
    Let $C$ be a class of simplicial complexes, then $|C|$ is closed under retracts if and only if $C$ is closed under bisimplicial images and triangulations. 
\end{corollary}
\begin{proof}
    $\PL-$retracts of $|\Sigma|$ are retracts of a triangulation of $\Sigma$ which are bisimplicial images by the lemma above.
\end{proof}

\subsection{Compact Goldblatt-Thomason for polyhedra}

Recall that the frames we defined are models for a modal logic (described extensively in \cite{DBLP}):

Essentially propositions are interpreted in unions of open subpolyhedra, $[\lozenge\phi]$ is the least closed subpolyhedron containing $[\phi]$.

We connect those as simplicial models by associating every simplex of a complex $\Sigma$ to its relative interior in $|\Sigma|$, thus obtaining a neighborhood-description:

\begin{corollary}
    If $\Sigma$ is a simplicial complex, for each point $x\in |\Sigma|$ there exists a unique face $\sigma\in\Sigma$ such that $x\in \sigma$, and 
    such $\Sigma$ can be interpreted as a poset through the face (or coface) relation.
\end{corollary}

We will quickly review the fact that the semantics of polyhedra (defined in \cite{Dekker}, \cite{DBLP}) is equivalent to the frame semantics for the face posets of simplicial complexes modulo subdivision.

Recall
\begin{definition}[Polyhedral semantics]
Let $\Ll$ be the basic modal language. We define a polyhedral semantics as an evaluation map onto the set $\Oo_P$\footnote{This will annoy algebraic geometers}
 of subpolyhedra of $P$ open in the weak topology (defined in section $1$.)
\[\nu: \Ll \to \Oo_P\]
Thus 
\begin{itemize}
    \item $P,x\models\varphi\iff x\in \nu(\varphi)$,
    \item $\nu(\varphi\vee\psi)=\nu(\varphi)\cup\nu(\psi)$,
    \item $\nu(\lnot\varphi)=\nu(\varphi)^c$,
    \item $\nu(\lozenge\varphi)=\overline{\nu(\varphi)}$.
\end{itemize}
\end{definition} 

To recollect the semantics in the literature, notice that 
\begin{enumerate}
    \item We can see $|-|$ as realizing open subpolyhedra: that is, if $\sigma\in \Sigma$ is a simplex,
     $|\sigma|$ is the open simplex of $|\Sigma|$ corresponding to it, this way 
        \[|\Sigma|=\bigsqcup_{\sigma\in\Sigma}|\sigma|\] is a disjoint union.
    This means that the closure of a $|\sigma|$ is the union of $|\sigma|$ with its boundary (i.e. all the subsimplices of $\sigma$), 
    in other words $\overline{|\sigma|}=|\downarrow\sigma|.$
    \item Using the classic correspondence, we have that the coheyting algebra $\Sub(P)$ 
    of closed subpolyehdra is indeed the codomain of evaluation of 
    the cointuitionistic\footnote{CoHeyting, paraconsistent} fragment of $\Ll$
    \[\set{\varphi\in \Ll:\varphi\Leftrightarrow \lozenge \phi}.\]
\end{enumerate}

Let's quickly review the correspondence: 
\begin{lemma}[Polyhedral-simplicial semantic equivalence]\hfill 

Let $\Sigma$ be a simplicial complex. Recall that for all $x\in |\Sigma|$ there exists a unique $\sigma_x\in \Sigma$ 
such that $x\in |\sigma_x|$.
Then, employing the kripke semantics on the face poset
\[|\Sigma|,x\models \varphi\iff [\Sigma],\sigma_x\models\varphi\]
\end{lemma}
\begin{proof}
A proof can be found in \cite{DBLP}.
\end{proof}
Moreover, 
\begin{lemma}
\[\set{\varphi\in \Ll: P\models\varphi}=\bigcap_{\Sigma:|\Sigma|=P}\set{\varphi\in\Ll:\Sigma\models\varphi}\]

i.e. the logic of a polyhedron is the intersection of the logics of its triangulations.
\end{lemma}
\begin{proof}
A proof can be found in \cite{Dekker}.
\end{proof}


Here we can finally express the main result of this paper:
\begin{theorem}[Compact GTT for Polyhedra]\hfill 

A class of compact polyhedra is modally definable if and only if it's closed under disjoint unions, closed subpolyhedra and retractions. 
\end{theorem}
\begin{proof}\hfill 

   Using the previous lemmas, we know that a class of polyhedra is modally definable if and only if the class of all of its triangulations is modally definable:

   Let $P$ be an element of $\set{P\in\Pol: P\models\varphi}$, then suppose $|\Sigma|=P$, we have that for all $x\in P, P,x\models\varphi$, 
   thus for all $\sigma_x\in \Sigma, \Sigma,\sigma_x\models\varphi$. Conversely, suppose every triangulation of $P$ models $\varphi$, then $P\models\varphi$ as well, since
   $\phi$ is in the intersection of the logics of all its triangulations.
   
   Thus, the class of polyhedra is modally definable if and only if the class of simplicial complexes formed by its triangulations (which is automatically closed under subdivision for Alexander's theorem)
   is modally definable, i.e. its face posets are modally definable.  
   
   The face posets are modally definable if and only if they're closed under disjoint unions, generated subframes and bounded morphic images,
   
   this is the case if and only if the class of polyhedra is closed under disjoint unions, subcomplexes and bisimplicial images,

   and this is the case if an only if the class of polyedra is closed under disjoint unions, closed subpolyhedra and retractions.
\end{proof}